\subsection{The construct is not patenting}\label{sec:patent}

Within the firms that do both, trademark vocabulary and patenting move
independently. On a panel of trademark owners name-matched to PatentsView
patent assignees \citep{PatentsView}, each axis is regressed on
$\log(1 + \mathrm{patents})$ with firms absorbed and errors clustered on
the firm, estimated on 94{,}181 firm-years across 17{,}506 firms under
theme scoring (93{,}551 and 17{,}415 under term scoring) --- 134{,}075
matched firm-years before dropping the firm-years that contribute no
within-firm variation. The outcome is standardized, so $\beta$ is the
shift in standard deviations of the axis per unit of log-patents; one
unit is roughly a tripling of the patent count, and a firm going from
none to twenty moves about three units.

\begin{center}\small
\begin{tabular}{llrrr}
\toprule
Text counted as & Axis & $\beta$ (sd) & $t$ & Between-firm $r$ \\
\midrule
Themes & atypicality & $+0.012$ & $+3.0$ & --- \\
Themes & lead        & $-0.058$ & $-9.4$ & $-0.020$ \\
Terms  & atypicality & $-0.055$ & $-11.5$ & --- \\
Terms  & lead        & $+0.012$ & $+2.0$ & --- \\
\bottomrule
\end{tabular}
\end{center}

\noindent The grid is close to a mirror: whichever axis carries the small
signal under one way of counting the text carries the opposite sign under
the other. The theme-scored \emph{lead} coefficient ($-0.058$) and the
term-scored \emph{atypicality} coefficient ($-0.055$) are nearly the same
number attached to different quantities. Only those two cells are
established; the mirror's other diagonal is within noise, and the
orthogonality conclusion rests on the magnitude bound --- every cell
under $0.06$ standard deviations --- rather than on any cell's sign.
Nothing here reaches a tenth of a standard deviation.

One objection is timing: a patent is applied for before its invention is
in public use and a trademark when the firm is ready to sell, so patents
would lead. Firms do not proceed that linearly \citep{Seip2018}, and no
timing offset from $t-2$ to $t+1$ isolates a direction ($-0.054$ to
$-0.089$ sd under theme scoring, $+0.012$ to $+0.024$ under terms);
among owners who both patent and raise a priced round, neither event
systematically precedes the other (median lag zero years).

The second objection is heterogeneity: a pooled near-zero could hide a
positive slope where the two rights are complements offsetting a negative
one where trademarks work alone. On a coarse, judgement-based sorting of
classes into complements and substitutes --- assigned in advance from the
class definitions, nothing estimated --- the slopes are $-0.000$ among
complements, $-0.026$ among substitutes, interaction $-0.008$
($t = -0.49$), and the complement classes do not agree in sign with one
another (cell detail in the supplementary material). Neither retention
supports the complementarity story, with the caveats that a coarse
sorting can miss a finer-grained one and that conditioning on a
successful assignee match pushes any interaction toward zero.

